\defmodule {seekl}

(To be rewritten....)

This program performs a search for the ``best'' (or ``worst'') multiple
recursive generators or multiply-with-carry generators of a given
form, based on one of the criteria $Q_T$ or $M_T$ defined in Section 1.3
It produces a report listing the retained generators, their properties,
and various statistics on the search.

The set of dimensions in which the test is applied can be partitioned
into a certain number of intervals, or categories,
and one can use a different selection criterion for the generators within
each category.  One can also impose bounds on the figure of merit
within each category.  See the data fields for $C$, {\it MinMerit}
and {\it MaxMerit\/} below.
For example, one can consider only the generators with $M_8 \ge 0.6$,
and among these, retain the list of generators with the {\em smallest\/}
value of $M_{12}$.
As another example, one can retain the 2 generators with the highest $M_8$,
the 2 generators with the highest $M_{16}$ and the eight generators
with the highest $M_{32}$.

One can search for combined MRGs with $J$ components, or simple MRGs $(J=1)$,
or multiply-with-carry (MWC) generators.
For a MWC, one simply analyzes the corresponding LCG, which is a special
case of an MRG.  Therefore, in what follows, we use the term
`MRG component' to denote either and MRG or a MWC.
For a simple MRG (or for each component, in case $J > 1$), with
given modulus $m$ and order $k$, the program searches for vectors
of multipliers inside the
region bounded by the vectors $b = (b_1,\dots,b_k)$ and $c = (c_1,\dots,c_k)$
such that $-m < b_i\le c_i < m$ for each $i$.
The search can be exhaustive in that region, or random.
One can search only among maximal period generators (for each component),
or not consider the period and examine only the lattice structure.
The former (checking maximal period conditions) can be done only
if $m$ is prime, or if $k=1$ and $m$ is a power of a prime.
The program can also list the retained generators in a file, in a format
more compact than for the result file, and can re-use that file as input
to the program, in a later run.
This could be useful, for example, if one wishes to perform first a screening
over a large region, based on a criterion that does not require expensive
computations, and then do a second pass over the retained generators,
based on a more stringent criterion, such as looking at the lattice structure in
higher dimensions, and/or verifying the results by performing all computations
using error bounds.


%%%%%%%%%%%%%%%%%%%%%%%%%%%%%%%%%%%%%%%%%%%%%%%%%%%%%%%%%%%%
\subsection* {Method of search}

For an exhaustive search for MRGs, all vectors of multipliers
of the form $a = (a_1,\dots,a_k)$
such that $b_i\le a_i\le c_i$ for $i=1,\dots,k$ will be examined, for
a total of $N_v = \prod_{i=1}^k (c_i-b_i+1)$ vectors.
This holds for each component.
Therefore, if there are $J$ components and $N_{v,j}$ vectors are
examined for component $j$, then a total of $\prod_{j=1}^J N_{v,j}$
generators are examined.

For a random search for MRGs, we fix a number of subregions (clusters)
we want to examine, and the size $h_i$ of each subregion in dimension $i$,
for $i=1,\dots,k$.  The program will examine a total of
$n\prod_{i=1}^k h_i$ vectors of multipliers (for each MRG component)
by repeating $n$ times the following:
% \begin {verse}
  For $i=1,\dots,k$, generate $\alpha_i$ randomly, uniformly over the set
  $\{b_i,\dots,c_i-h_i+1\}$;  then, examine all the vectors
  $a = (a_1,\dots,a_k)$ such that $\alpha_i\le a_i\le \alpha_i+h_i-1$
  for each $i$.
% \end {verse}

When examining a vector $a$, the program first checks if the maximal
period conditions are satisfied, if this is required.
For prime modulus $m$, the condition (a) (see Section 1.2) is verified
only once for each distinct value of $a_k$ (which corresponds to
$\prod_{i=1}^{k-1} h_i$ different vectors).
To verify the maximal period conditions, the factorizations of $m-1$ and
$r = (m^k-1)/(m-1)$ are required.  They can be found by the program,
if desired, or provided by the user in a file (see below).
We recall that factoring $r$ can take {\em huge} amounts of time.
So, avoid redoing the factorization unnecessarily.
For the MRGs, these factorizations are necessary only when $m$ is
prime and maximal period is required.

If $a$ is not rejected by the maximal period test, then we move forward
to the next MRG component and try all the vectors for that next component
(by exhaustive or random search) and examine their combination with
the currently examined multipliers for the previous components.
For each combined generator, the values of $d_t$ and/or $q_t$ are 
computed for dimensions $k+1,\dots,T$.  

The program always keeps lower and upper bounds
on the figure of merit ($M_T$ or $Q_T$), in each dimension,
for the generator to be worth considering.  
The initial values of these bounds are given by the user
in the fields {\em MinMerit} and {\em MaxMerit}.
The lower bound can be 0.0 and the upper bound can be 1.0, 
which means that there can be no effective bounds for some
categories if desired.

As soon as a generator has a figure of merit below the lower
bound in a given dimension, or above the upper bound for its
category (after the computations for all the dimensions in this
category have been completed), then this generator is immediately
discarded and no further computations are made for it.
This can save enormous amounts of time in the case of very large
searches up to high dimensions, because with good bounds,
few generators will reach the large dimensions.

During execution, only the bounds for the last category can be
modified.  If the figure of merit for the last category is to be
{\em maximized\/}, when we have found enough (i.e., {\em NbGen($C$)})
generators with a figure of merit $\ge \sigma$, where $\sigma$
is larger than the lower bound for the last category, then we raise
this lower bound to $\sigma$.
Similarly, if we minimize in the last category, we can lower 
the upper bound when we have enough generators beating the bound.
In the case where the figure of merit is to be {\em maximized in 
all\/} the categories, then a generator is also discarded
as soon as its figure of merit in {\em any\/} dimension gets below 
the lower bound of the last category.

In the case where the selection criterion is {\tt Spectral} or
{\tt SpectralH}, and if $T>8$, the distances between
hyperplanes are computed for dimensions up to $T$, but the selection of
generators is based only on $M_{8}$, because the Hermite constants
defining $d_t^*$ in that case are known only for $t\le 8$.

The execution (cpu) time is checked before testing each new generator.
When it exceeds the cpu time limit given in the data file,
the search is aborted and the partial results are printed.

%%  DATA FILE  %%%%%%%%%%%%%%%%%%%%%%%%%%%%%%%%%%%%%%%%
\setbox0=\vbox {\hsize=12.0cm \it \begin {tabular}{l}
           ReadGenFile ~$[$\<file0\>$]$ \\
           $J$ \\
           $\left.\begin {array}{l}
           \mbox {TypeGen} \\
           m \\
           k \\
           w \\
           \mbox {PerMax} \\
           \mbox {ImplemCond} ~[\mbox{\<NbMaxBits ~HighestBit\>}] \\
           F1 ~[\mbox{\<file1\>}] \\
           F2 ~[\mbox{\<file2\>}] \\
           b_1 \\
           c_1 \\
           \vdots \\
           b_k \\
           c_k \\
           \mbox {SearchMethod $[$n H $H_k]$ } \\
           \end {array} \ \ \ \right\} \mbox { (Repeat $J$ times) }$ \\
           $C$ \\
           NbGen$(1)$ MinMerit$(1)$ MaxMerit$(1)$ 
                $d_1$ $t_{11}$ $\cdots$ $t_{1,d_1}$ \\
           \vdots \\
           NbGen$(C)$ MinMerit$(C)$ MaxMerit$(C)$ 
                $d_C$ $t_{C1}$ $\cdots$ $t_{C,d_C}$ \\
           Criterion  ~\<{\em Invert\/}\> ~\<{\em Norm\/}\> \\
           LatticeInfo \\
           LatticeType \\
           LaGroupSizes ~Spacing\\
           VerifyBB \\
           MaxNodesBB \\
           TimeLimit \\
           $S1 ~S2$ \\
           ResultForm \\
\end {tabular}
}
\begin{figure}[htb] \centerline {\boxit{\box0}}
\caption { Data file format for {\tt seekl}. }  \label {seekl.dat}
\end{figure}

%%%%%%%%%%%%%%%%%%%%%%%%%%%%%%%%%%%%%%%%%%%%%%%%%%%%
\subsection* {The data file}

The data for {\tt seekl} must be placed in a file with extension
``{\tt .dat}'', according to the format displayed on Figure~\ref{seekl.dat}.
The fields in square brackets are optional (depending on the value taken
by the first field on the line).
The meaning of all data fields is explained below.
To run the program, type ``{\tt seekl \<\em file}\>'', where \<{\em file}\> is
the name of the data file, without extension.
The results will be in a file with the same name, with extension
``{\tt .res}'' or ``{\tt .gen}'' (see {\em ResultForm\/}).

Comments may be inserted after data, on the same line,
separated from the data by at least one blank. 
Moreover, any line starting with ``\%'' or ``\#'' is considered
as a comment.

The values of $m$, $b_i$, and $c_i$ in the data file can be given in one
of the two following formats:
\vspace {-10pt}
\begin {itemize}
\item [a)]  An integer giving the value directly, in base 10.
  In this case, there {\em must\/} be some other non-numeric text 
  (e.g., a comment) on this data line after the integer.
\item [b)]  Three integers $x$, $y$, and $z$ on the same line, separated
  by at least one blank.  The retained value will be $x^y+z$ if $x\ge0$, and
  $-((-x)^y+z)$ if $x<0$.  The value of $y$ must be positive.
  For example, $(x\ y\ z) = (2\ 5\ -1)$ will give 31, while
  $(x\ y\ z) = (-2\ 5\ -1)$ will give $-31$ (not $-33$).
\end {itemize}
For the program {\tt seekl}, all these numbers must fit in a {\tt LONGINT}.
For larger numbers, one must use the program {\tt seeks}.

%%%%%%%%%%%%%%%%%%%%%%%%%%%%%%%%%%%%%%%%%%%%%%%%%%%%%%%%%%%%%%%
\subsection* {Meaning of the data fields}

\begin {description}
%
\item [{\em ReadGenFile} and \<{\em file0}\> :]
  {\tt BOOLEAN} and file name (without extension).
  When {\em ReadGenFile} is {\tt FALSE},
  the search is made according to the values of the fields below.
  When {\tt TRUE}, the generators to be looked at are those
  listed in the file \<{\em file0}\>{\tt .gen}.
  This must be a file of type ``{\tt .gen}'', produced by this program
  with the {\tt GEN} option for the {\em ResultForm} data field.
  In that case, only those generators listed in that file are examined
  and the vectors $b$ and $c$ below are not used.
%
\item [{\em J} :]
  Number of components in the combined generators.
  Must be an integer in the range $[1..{\tt MaxJ}]$.
  When $J > 1$, we look for combined generators.
%
\item [{\em TypeGen} :] Can be:
  \begin {description}
  \item [{\tt MRG}] means that this component is an MRG.
  \item [{\tt MWC}] means that this component is a multiply-with-carry 
    (MWC) generator.
    Each MWC generator is converted by the program to its
    corresponding LCG (see, e.g., \cite{rCOU97a,rLEC98a}).
  \item [{\tt MMRG}] means a matrix MRG.
  \end {description}
%
\item [{\em m, k}, \<{\em w}\> :]
  Modulus, order of the recurrence, and size of the matrix-type 
  coefficients.
  Must be positive integers, with $k<$ {\tt MMaxDim}.
  For $m^k >= 2^{31}$, use the program {\tt seeks} instead of {\tt seekl}.
  The value of $w$ should be specified only for MMRG generators.
  Otherwise it is assumed to be 1.
%
\item [{\em Permax} :]
  {\tt BOOLEAN} variable.  
  {\tt TRUE} if maximal period is required, {\tt FALSE} otherwise.
  When set to {\tt TRUE}, $m$ must be expressed in the data file in the
  form (b): $(x\ y\ z)$, otherwise {\em Permax\/} will be put back to
  {\tt FALSE}.  The software assumes that $m$ is prime, 
  unless $z=0$ and $y > 1$, in which case it assumes that $x$ is prime.
  In the latter case, one must have $k=1$, otherwise {\em Permax\/}
  will be set back to {\tt FALSE}.
%
\item [{\em ImplemCond} and \<{\em NbPow2 ~HighestBit}\> :]
  Can be {\tt NoCond}, {\tt AppFact}, or {\tt SumPow2}.
  If {\tt NoCond}, then no conditions are imposed on the
  multipliers $a_i$.
  If {\tt AppFact}, then the multipliers must satisfy the 
  ``approximate factoring'' condition $a_i (m\mod a_i) < m$ for each $i$.
  MRGs are usually easier to implement under this condition \cite{rLEC90a}.
  If {\tt SumPow2}, then the positive integers {\em NbPow2}
  and {\em HighestBit\/} must appear and they indicate that 
  for each $i$, the multiplier $a_i$ must be the sum of at most 
  {\em NbPow2} (positive or negative) powers of 2, 
  with the highest power of 2 not exceeding $2^e$ in absolute value,
  where $e =$ {\em HighestBit}.  
  For example, if {\em SumPow2\/} = 2 and {\em HighestBit\/} = 30,
  there are $30 \times 31 /2$ possibilities for choosing the 2 powers of 
  2 and 4 possibilities for choosing their signs, yielding
  1860 cases where $a_i$ is obtained from exactly 2 powers of 2.
  If one adds the $62$ cases where $a_i$ is $\pm$ a power of 2,
  this gives a total of 1922 possibilities for $a_i$.
%  In the case of {\tt SumPow2}, the bounds
%  $b_j$ and $c_j$ on the multipliers (see below) are automatically
%  (implicitly) reset to $b_j = -(2^{e+1}-1)$ and $c_j = 2^{e+1}-1$.
%
\item [$F1$ and \<{\em file1}\> :]
  This line of data (and also the following one) is used only if 
  maximal period is required and $m$ is assumed to be prime
  (see the {\em Permax\/} field).  Otherwise, the program just skips it
  (but the line must be there).
  $F1$ indicates how the factors of $m-1$ are to be found and
  \<{\em file1}\> is a file name.
  The values allowed for $F1$ are {\tt (Decomp, DecompWrite, Read)}.
  \begin {description}
  \item  [{\tt Decomp}] means that the program itself will factorize
    $m-1$.  In this case, the field \<{\em file1}\> is not used and can be
    omitted.
    To factorize, the program uses the procedure
    {\tt Factorize} in the module {\tt SUPFACT} of the package SENTIERS
    \cite {iLEC88d}, with no CPU time limit.
    It is the responsibility of the user to make sure that the
    factorization will take a reasonable time.
  \item [{\tt DecompWrite}] means the same as {\tt Decomp}, except that
    the program will also write the factors (larger than 1) in the file
    \<{\em file1}\>, one factor per line.
  \item [{\tt Read}] indicates that $m-1$ is already factorized and that
    the factors will be read from file \<{\em file1}\>, in the same format.
    The factors need not be sorted, but must be one per line, and repeated
    factors must be on successive lines.  The factorization must be
    complete and the program will check if the product of all the factors
    is really equal to $m-1$.
  \end {description}
%
\item [$F2$ and \<{\em file2}\> :]
  This data line is similar to the previous one, except that it concerns
  $r=(m^{kw}-1)/(m-1)$ instead of $m-1$.
  In this case, it is possible that $r$ be prime when $m$ is prime
  (in contrast to $m-1$, which is then even).  
  Therefore, the additional value {\tt Prime} is
  allowed for $F2$, so that the set of possible values is
  {\tt (Decomp, DecompWrite, Read, Prime)}.
  The module {\tt SUPFACT} in SENTIERS \cite{iLEC88d} can be used independently
  to factorize $r$ or to check its primality.
  \begin {description}
  \item [{\tt Prime}] indicates that $r$ is prime.
  \end {description}
%
\item [$b=(b_1,\ldots,b_k)$ \rm and $c=(c_1,\ldots,c_k)$~:]
  The $b_i$ and $c_i$ are integers such that
  $-m < b_i \le c_i < m$ for $i=1,\dots,k$.
  They determine the boundary of the (rectangular) area of search.
%  When {\em ImplemCond} = {\tt SumPow2}, these values are not used.
  For MMRGs ($w>1$), each $b_i$ and $c_i$ is replaced by a matrix of
  size $w\times w$.  Each of these matrices is given line by line,
  $w$ entries on each line of the data file.
%
\item [{\em SearchMethod} and $n,\, H,\, H_k$ :]
  {\em SearchMethod} can be {\tt Exhaust} or {\tt Random}.
  \begin {description}
  \item  [{\tt Exhaust}] means that the search will be exhaustive over
    all the region determined by $b$ and $c$.
    The other parameters of this line are then useless and can be omitted.
  \item  [{\tt Random}] asks for a random search, performed as described
    in the first subsection of the description of {\tt seekl}.
    The integer $n$ gives the number of subregions (clusters) to examine. 
    If $w=1$, $H$ determines the size of these subregions, except for
    the $k$th element of the vector $h$, where $H_k$ determines the size.
%    The parameter $Hk$ is not used if $k = 1$.
    The vector $h=(h_1,\ldots,h_k)$ is computed by the program as follows:
    $h_i = \min (H, c_i-b_i+1),\ i=1,\ldots,k-1$, and
    $h_k = \min (H_k, c_k-b_k+1)$.
    If $w > 1$, ?????
  \end {description}
%
\item [$C$:] 
  The program will retain $C$ lists (or categories) of generators, 
  according to the specifications given below.
  One must have $1 \le C \le {\tt MaxCat}$.
%
\item [{\em NbGen$(1)$ ~MinMerit$(1)$ ~MaxMerit$(1)$} 
        $d_1$ $t_{11}$ $\cdots$ $t_{1,d_1}$]
\item [ \hskip 1in\vdots ] \
\item [{\em NbGen$(C)$ ~MinMerit$(C)$ ~MaxMerit$(C)$} 
        $d_C$ $t_{C1}$ $\cdots$ $t_{C,d_C}$] 
  \ \\ \ \\
  Integers which must satisfy the constraints: 
  \begin {eqnarray*}
   & 0\le {\em NbGen}(i) \le {\tt MaxNbGen},\\
   & 0.0 \le {\it MinMerit}(c) \le {\it MaxMerit}(c) \le 1.0 &
            \qquad\mbox{ for all } c, \\
   & {\it MinMerit}(c) \le {\it MaxMerit}(c-1) & \qquad\mbox{ for } c > 1,\\
   & d_1 \le \cdots\le d_C, \\
   & {\tt BMaxDim} \ge t_{c1} \ge \cdots\ge t_{c d_c}\ge d_c &
            \qquad\mbox{ for all } c, \\ 
% t_{1d} \le \cdots\le t_{C d} \mbox{ for all } d. 
%  -> Try keeping this constraint out.
  \end {eqnarray*}
  where {\tt BMaxDim} is defined in the module {\tt CONFIG}.

  There will be {\em NbGen$(c)$} generators in the $c$th list.
  The real numbers {\em MinMerit$(c)$ and MaxMerit$(c)$} represent 
  the minimal and maximal values of the figure of
  merit to keep a generator in the $c$th list.
  That is, only the generators that satisfy 
 \eq  
  \mbox{\it MinMerit}(c) \le\sigma\le \mbox{\it MaxMerit}(c) 
                                                   \label{con-merit}
 \endeq 
  are considered for category $c$ and for the categories $c' > c$.  
  Because of this, {\it MinMerit}$(c)$ must be nondecreasing in $c$,
  because the lower bound for category $c$ also
  apply categories $c' > c$.
  As soon as a generator does not satisfy the constraint (\ref{con-merit})
  in a given dimension for a given category $c$,
  it is discarded and no more time is spent to test it for $c' > c$.
  When looking for good generators, one normally sets 
  {\em MaxMerit}$(c)$ to 1.0 for each $c$.

  For category $c$, the test (Beyer, spectral, \dots, depending on
  the criterion) is performed in dimensions $kw+1$ to $t_{c1}$ 
  for the vectors of successive values,
  the pairs $(u_1, u_{j})$ are tested for $1 <j\le t_{c2}$,
  the triples $(u_1, u_{j}, u_{\ell})$ are tested for $1 <j <\ell\le t_{c3}$,
  and so on.  
  If lacunary indices are specified in the field {\em LaGroupSizes},
  all of this is applied to the vectors obtained {\em after\/} taking
  the lacunary indices, and the test for successive values starts in 
  2 dimensions instead of $kw+1$.
  To {\em minimize\/} (instead of maximizing) the criterion within a 
  given category $c$, place a {\em negative sign\/} in front of the value
  of $d_c$ in the data file.
  The program will then retain the {\em worst\/} instead of the
  {\em best\/} generators with respect to the category $c$.
%
\item [{\em Criterion} ~\<{\em Invert\/}\> ~\<{\em Norm\/}\> :]
  Specifies the merit criterion for ranking the generators for each category.
  The field {\em Invert\/} can be either {\tt Invert} or left blank.
  In the former case, the value of the criterion is
  replaced by its multiplicative inverse in the results.
  The admissible values of {\em Criterion\/} are:
  \begin {description}
  \item [{\tt Beyer}] means that the criterion is $Q_T$.
    The program will retain the generators with the largest 
    (or smallest) $Q_T$ in each category.
  \item [{\tt Spectral}] means that the criterion is $M_T$.
    The program will retain the generators with the largest values
    (or smallest) $M_T$ in each category.
  \item [{\tt SpectralP}] is similar to {\tt Spectral}, except that
    it is based on the shortest vector in the {\em primal\/} lattice.
  \item [{\tt SpectralL1}] is similar to {\tt Spectral}, except that
    the length of the (dual) vectors are measured with the $L_1$ norm,
    minus 1.  The length of the shortest dual vector is then an upper
    bound on the minimal number of hyperplanes that cover all the points
    (\cite{rDIE75a} and \cite{rKNU97a}, Exercises 3.3.4-15 and 16).
  \item [{\tt Palpha} {\em Alpha}] means that the criterion is $P_\alpha$,
    with the value of $\alpha$ specified just after the keyword {\tt Palpha}.
    Specialized (much faster) implementations have been made for $\alpha = 2, 4, 6, 8$.
  \end {description}
  If {\em Criterion\/} is {\tt Spectral} or {\tt SpectralP}, 
  then {\em Norm\/}
  must appear and it indicates which type of normalization is used
  in the definition of $S_t$.  The admissible values are
  \begin {description}
  \item [{\tt Hermite}] means that we use Hermite's bound.
    In this case, $S_t$ is computed only up to dimension 8
    and the criterion is $M_{\min(T, 8)}$.  
    This is because for $t>8$, $d_t^*$ is
    unknown, so $M_t$ cannot be computed.
  \item  [{\tt Laminated}] means that we use for $d_t^*$ the value
    of $d_t$ that corresponds to the laminated lattice in dimension $t$.
    The criterion is then $M_{\min(T, 48)}$ 
    (the programs knows this $d_t^*$ only for $t\le 48$).
  \item  [{\tt Rogers}] means that $d_t^*$ is obtained from Rogers'
    bound.  The criterion is then $M_{\min(T, 48)}$.
  \end {description}
  If {\em Criterion\/} is {\tt Palpha}, the field {\em Norm} must also
  appear.  The admissible values for {\tt Palpha} are:
  \begin {description}
%  \item [{\tt ????? }] (Check the paper with Hickernell) ?????
  \item [{\tt Weights}] followed on the next line by positive real numbers 
  $\beta_0,\dots,\beta_t$ where $t= t_{C,1}$.  This means that the criterion
  will be $\tilde P_{\alpha}$ defined in (\ref{eq:Palpha}).
  In this case, $\alpha$ must be even.  The $\beta_j$'s are the weights
  given to the different dimensions in the weighted version of $P_{\alpha}$.
  \end {description}
%
\item [{\em LatticeInfo} :]
  Indicates which figures of merit we want to compute and print.
  The possibilities are the same as for the criterion, but we can compute
  and print more than one.  It suffices to write more than one keywords
  on this line.
  Note that $q_t$ typically requires much more time to compute than $d_t$.
  \begin {description}
  \item [{\tt Beyer} :] Prints the values of $q_t$ up to dimension $T$.
  \item [{\tt Spectral} :] Prints the values of $d_t$ up to dimension $T$ and
    $S_t$ for $t\le\min (T, T^*)$, where $T^*=8$ if {\em Norm} = 
    {\tt Hermite} and $T^* = 48$ otherwise.
  \item [{\tt SpectralP}] is similar to {\tt Spectral}, except that
    it is based on the shortest vector in the {\em primal\/} lattice.
  \item [{\tt SpectralL1}] is similar to {\tt Spectral}, except that
    the length of the (dual) vectors are measured with the $L_1$ norm,
    minus 1.  The normalization is based on Minkowski's bound
    (\cite{sFIS96a}, p. 618).
  \item [{\tt Palpha} {\em Alpha} :] prints $P_\alpha$
    with the value of $\alpha$ specified.
  \end {description}
%
\item [{\em LatticeType} :]
  Indicates whether to analyze the lattice generated by all possible states,
  or a sublattice generated by the set of recurrent states or by a
  subcycle of the generator.
  The admissible values are {\tt (Full, Recurrent, Orbit, PrimePower)}.
  \begin {description}
  \item [{\tt Full} :] The complete lattice, generated by all possible
     initial states, will be analyzed.
  \item [{\tt Recurrent} :] If the (combined) generator has transient
     states, then the lattice analyzed will be the sublattice generated by
     the set of recurrent states.
  \item[{\tt Orbit} :] The grid generated by the (forward) orbit of 
     a state of the (combined) generator is analyzed. This state
     is specified as follows. 
     On the following $J$ lines, the initial state for each
     component must be given. This is an integer vector with a number 
     of components equal to the order of the component.
  \item [{\tt PrimePower} :] In the case where some component is an MLCG
     whose modulus is a power of a prime $p$, then the states visited over
     a single orbit (subcycle) of that component generate a sublattice
     (when $a \equiv 1 \pmod{p}$) or belong to the union of
     $p-1$ sublattices (otherwise).  If {\em LatticeType\/} takes this value,
     if a component is an MLCG ($k=1$), 
     and if the modulus of that MLCG is given
     in the data file in the form (b): $(x\ y\ z)$ with $z=0$ and $x$ prime, 
     then what is analyzed
     is one of those sublattices.  This is done by dividing the modulus
     by the appropriate power of $p$, as described in \cite{rLEC97c}.
     For example, if $p=2$ and $a \mod 8 = 5$, then the modulus is divided
     by 4 as in \cite{rFIS90a,rKNU81a}.
  \end {description}
%
\item [{\em LaGroupSizes, ~Spacing\/} :]
  These data fields are positive integers, used to introduce lacunary indices.
  If the respective values are $s$ and $d$, then we will analyze the
  lattice structure of vectors of the form
  $(u_{i+1}, \dots, u_{i+s}, u_{i+d+1},\dots, u_{i+d+s}, u_{i+2d+1},\dots,
  u_{i+2d+s}, \dots)$, formed by groups of $s$ successive values,
  taken $d$ values apart.
  To analyze vectors of successive values (as usual), take $s=d=1$
  or $s$ larger or equal to {\em MaxDim}.
  To analyze lacunary indices that are not evenly spaced,
  put $s = -t$ where {\em $t=$MaxDim\/} and then, on the $t$ lines that
  follow, give the $t$ lacunary indices $i_1,\dots,i_t$, which are to
  be interpreted as in Section 1.4.
%
\item [{\em VerifyBB} :]
  Indicates whether or not the results of the branch-and-bound procedures
  will be verified with formal error bounds on the numerical computational
  errors.  The verification (if any) is performed using either
  {\tt REDBAS2} or {\tt REDBAS3}.
  When such a verification is not performed, the results are not
  formally guaranteed, but computations are performed an order of
  magnitude faster.  The possible values are:
 \begin {description}
 \item [{\tt VerifyAll2} :] Performs the verification for all the generators
   examined, using {\tt REDBAS2}.
 \item [{\tt VerifyRetained2} :] Performs the verification only for the
   {\em NbGen\/} ``best'' generators retained, using {\tt REDBAS2}.
 \item [{\tt VerifyAll3} :] Performs the verification for all the generators
   examined, using {\tt REDBAS3}.
 \item [{\tt VerifyRetained3} :] Performs the verification only for the
   {\em NbGen\/} ``best'' generators retained, using {\tt REDBAS3}.
 \item [{\tt NoVerify} :] Performs no verification at all.
 \end {description}
%
\item [{\em MaxNodesBB} :]
  An integer giving the maximum number of nodes to be examined in any
  given branch-and-bound procedure when computing $d_t$ or $q_t$.
  When that value is exceeded, the branch-and-bound is stopped and the
  generator is rejected.  The number of generators rejected for that
  reason is given in the results.
  A small value of {\em MaxNodesBB} will make the program run faster
  (sometimes much faster), permitting to examine more generators,
  but will increase the chances of rejecting good generators.
%
\item [{\em TimeLimit} :]
  A real number giving the maximal CPU time, in hours, given to the program
  for performing its search.
  When this time-limit expires, the partial results are printed, with a message.
  For a random search, the number of subregions printed includes
  the last region searched (whose search may not be finished).
%
\item [$S1,\, S2$ :]
  Those integers are the seeds of the generator used during the
  search, if necessary; e.g., to generate the $\alpha_i$'s.
  Before starting the search, the program calls
  {\tt SetSeed($S1, S2$)} form the module {\tt SUPRAND} of SENTIERS
  and the module {\tt RAND} of SIMOD.
  To perform a different random search in a region already studied,
  just change these seeds.
  One must have $1\le S1\le 2147483562$ and $1\le S2\le 2147483398$.
%
\item [{\em ResultForm} :]
  Selects in which form the results will be given.
  The possible values are {\tt (Terminal, RES, GEN, RESGEN)}.
  Lowercases are also allowed.
  \begin {description}
   \item [{\tt Terminal}] indicates that the results will appear only on the
     terminal screen.
   \item [{\tt RES}] says that the results will be in a file with the same
     name as the data file, but with extension ``{\tt .res}''.
   \item [{\tt GEN}] says that the retained generators will be listed in
     a file with the same name as the data file, with extension ``{\tt .gen}''.
     This file can then be taken as input to the same program, for
     example to perform a second pass with a more stringent criterion
     or to compute higher dimensional lattice ``measures'' for the
     retained generators.
   \item [{\tt RESGEN}] means that the results will be in a ``{\tt .res}''
     file {\em and} the retained generators will be listed in a ``{\tt .gen}''
     file, as with {\tt RES} and {\tt GEN} above.
  \end {description}
\end {description}

